% =========================================================================
% NPL finite-domain quantified paper — release v0.7.1 (2026-08-16).
% Release normalization: stable Lean/mathlib pair, self-contained replay definition,
% refreshed related work, synchronized artifact and reproducibility metadata.
% v0.7.1: citation corrections; public repository and archival DOI synchronization.
% v0.7.0: full flat replay compiler and universal membership-certificate bridge proved.
% v0.6: membership-selecting replay repair, soundness and regression boundary;
% universal completeness of the new certificate recorded as conjectural (superseded).
% v0.5: admissible replay-certificate bridge refuted by a nonempty fold cascade.
% v0.4: genuine signature-indexed model semantics, raw/signature bridge,
% and completeness against signature consequence.
% v0.3: fixed-signature/arity layer and well-formed completeness scope added;
% constructor-replay status corrected.
% Companion to papers/npl-core. Headline: semantic completeness of the
% macro-free core extensional tableau (docs Thm 3.74), plus Thm 3.75.
% Manuscript numbering is presentational; the stable numbering lives in
% docs/ (project rule R3). Results marked [L] name sorry-free Lean
% declarations in the Nullivance artifact.
% Build: pdflatex main && bibtex main && pdflatex main && pdflatex main
% =========================================================================
\documentclass[11pt]{article}

\usepackage[T1]{fontenc}
\usepackage[utf8]{inputenc}
\usepackage{lmodern}
\usepackage{amsmath,amssymb,amsthm}
\usepackage{booktabs,array}
\usepackage[colorlinks=true,linkcolor=blue,citecolor=blue,urlcolor=blue]{hyperref}
\hypersetup{
  pdftitle={Finite-Domain Quantified Nullivance Logic: A Machine-Checked Completeness Theorem for a Macro-Free Extensional Tableau},
  pdfauthor={Trinh Tung Lam},
  pdfsubject={A formally verified finite-domain quantified extension of nullivance logic},
  pdfkeywords={nullivance logic, quantified logic, paraconsistent logic, tableaux, completeness, Lean 4}
}
\usepackage[capitalize]{cleveref}
\usepackage{microtype}
\newcommand{\lb}{\discretionary{}{}{}}
\emergencystretch=3em

\newtheorem{theorem}{Theorem}[section]
\newtheorem{lemma}[theorem]{Lemma}
\newtheorem{proposition}[theorem]{Proposition}
\newtheorem{corollary}[theorem]{Corollary}
\theoremstyle{definition}
\newtheorem{definition}[theorem]{Definition}
\newtheorem{remark}[theorem]{Remark}
\newtheorem{example}[theorem]{Example}

% ----- notation (shared with the companion paper) -----
\newcommand{\FOUR}{\mathbf{4}}
\newcommand{\vT}{\mathbf{T}}
\newcommand{\vF}{\mathbf{F}}
\newcommand{\vB}{\mathbf{B}}
\newcommand{\vN}{\mathbf{N}}
\newcommand{\hmz}{\oplus}
\newcommand{\proj}[1]{\pi_{#1}}
\newcommand{\ind}[1]{\mathbf{1}[#1]}
\newcommand{\Tp}{T^{+}}\newcommand{\Tm}{T^{-}}
\newcommand{\Fp}{F^{+}}\newcommand{\Fm}{F^{-}}
\newcommand{\lean}[1]{\texttt{#1}}
\newenvironment{leanrefs}
  {\par\smallskip\noindent\begin{minipage}{\linewidth}\small\raggedright}
  {\end{minipage}\par\smallskip}
% finite-FO specific
\newcommand{\qder}{\vdash_{\mathrm{core}}}
\newcommand{\qmod}{\models_{Q,n}}
\newcommand{\grd}{\mathrm{gr}}
\newcommand{\qsize}{\mathrm{qsize}}

\title{Finite-Domain Quantified Nullivance Logic:\\
A Machine-Checked Completeness Theorem\\
for a Macro-Free Extensional Tableau\thanks{%
Release v0.7.1 (August 2026). Companion to \emph{Nullivance Propositional Logic:
Threshold-Signed Consequence on the Unit Square with a Machine-Checked
Completeness Theorem}. Results
marked [L] are verified sorry-free in the Lean~4 artifact described in
Appendix~\ref{app:lean}. Related-work positioning is backed by a dated,
source-by-source literature record supplied with the artifact; novelty claims
are search-scoped.}}
\author{Trinh Tung Lam\thanks{Contact: \texttt{lamtung0481@gmail.com}.}}
\date{August 2026 --- release v0.7.1}

\begin{document}
\maketitle

\begin{abstract}
We extend nullivance propositional logic (NPL) --- a paraconsistent,
two-channel logic on the unit square with a harmonization connective $\hmz$
and four threshold signs $\Tp,\Tm,\Fp,\Fm$ --- by finite-domain quantifiers.
The language is function-free first-order over a fixed signature with crisp equality; models
interpret predicates over a fixed finite nonempty domain by values in the
four-element matrix $\FOUR$, and the finite universal (existential)
quantifier takes the pointwise minimum (maximum) on the truth channel and
the dual on the falsity channel. Main results, each machine-checked in
Lean~4: (i)~the exact threshold projection of the propositional theory
lifts to the quantified language over finite domains; (ii)~four tempting
proof routes are \emph{refuted} by explicit countermodels: the naive
lifted tableau (crisp $x{=}x$ is valid but underivable), the bare
grounding reduction to propositional consequence, and a dispatcher
shortcut and the head-sensitive certificate bridge in the simulation
program; (iii)~the
repaired \emph{core extensional tableau}, whose closure conditions add
four crisp-equality clauses and two ground-closure clauses to the
decomposition rules, is sound and \emph{complete}: a branch closes exactly
when no finite model satisfies it. The completeness proof is a quantified
tableau engine driven by a domain-weighted measure
($\qsize(\forall x\,\varphi)=(n{+}1)\cdot\qsize(\varphi)+1$ over a domain
of size $n{+}1$) under which every rule strictly decreases branch weight,
together with a canonical-model construction at the literal stage keyed by
an injective ground-atom encoding. As a corollary, closure of the grounded
propositional branch under rigid constraints reaches the core calculus with
no propositional-simulation macro --- settling the semantic bridge. We
also document the failure analysis that led from a derivation-by-derivation
simulation program to the weighted-measure proof, and a conservative
membership-selecting repair of the replay certificate. The repair is
sound, passes the motivating regression, and is complete for every
admissible ground-closed trace. Its completeness is proved by fold-first
membership saturation followed by a domain-weighted compiler, rather than
inferred from semantic completeness.
\end{abstract}

\section{Introduction}\label{sec:intro}

The companion paper develops the propositional core of nullivance logic
(NPL): formulas valued in $[0,1]^2$ with independent truth and falsity
channels, a threshold $\tau\in(0,1]$ inducing four signs, an exact
projection onto the four-valued matrix $\FOUR$, and a four-signed analytic
tableau calculus with a machine-checked completeness theorem. This paper
asks the next question: what does a \emph{quantified} NPL look like, and
can its metatheory be machine-checked to the same standard?

We deliberately restrict to \textbf{finite nonempty domains}. The
restriction is not cosmetic: the exact-projection theorem — the load-bearing
bridge between the continuous semantics and the finite matrix — uses that
finite infima and suprema are attained, and it can fail over infinite
domains when a supremum is approached but not attained
(\cref{rem:infinite}). Finite domains keep the entire two-layer
architecture of NPL intact, and they are the natural first-order setting
for the intended applications of a threshold logic (finite data, finite
scenarios). Equality is \emph{crisp}: $x{=}y$ evaluates to $\vT$ or $\vF$
by identity of domain elements, never to a glut or a gap.

\paragraph{Contributions.} All results are verified sorry-free in Lean~4
unless explicitly marked otherwise; [L] markers name the verifying
declarations.

\begin{enumerate}
\item \emph{Finite-domain quantified NPL} (\cref{sec:syntax}): function-free
  syntax with predicate atoms, crisp equality, the NPL connectives
  including $\hmz$, and quantifiers evaluated by finite channel-wise
  extrema; quantifier duality under channel-swap negation; and the exact
  threshold projection for the quantified language.
\item \emph{A catalogue of machine-checked refutations}
  (\cref{sec:negative}): the naive signed tableau is incomplete (crisp
  reflexivity is valid but underivable); the bare grounding bridge to
  propositional NPL is unsound as a completeness route; and a natural
  dispatcher shortcut for the completeness proof fails on a two-element
  counterexample. Each failure shaped the final calculus.
\item \emph{The core extensional tableau and its exact semantic
  characterization} (\cref{sec:core,sec:completeness}): the headline
  theorem states that a finite branch closes in the macro-free core
  calculus \emph{iff} it has no finite model, hence derivability coincides
  with finite semantic consequence. The proof introduces a domain-weighted
  formula measure under which every rule of the calculus — including the
  $(n{+}1)$-ary quantifier rules — strictly decreases the branch weight,
  and closes the literal stage by a canonical model read off the branch
  through an injective ground-atom encoding.
\item \emph{The constrained-grounding corollary} (\cref{sec:grounding}):
  closure of the grounded propositional branch augmented by \emph{rigid
  constraints} (fixing the identity and equality atoms) yields core
  closure. Unlike the first completeness theorem of the artifact, no
  propositional-simulation macro rule appears anywhere: the propositional
  derivation is consumed through its soundness and the core closure is
  rebuilt semantically.
\item \emph{A documented failure and conservative repair} (\cref{sec:simulation}): the
  derivation-by-derivation simulation program (replay traces with residual
  fold tails) is developed to the point where its exact obstruction is
  isolated — the binary propositional decomposition cascade cannot be
  reassembled into the core's $(n{+}1)$-ary quantifier rules — and the
  weighted-measure proof is shown to dissolve the obstruction. A minimal
  admissible two-element fold cascade now refutes the proposed universal
  replay-certificate bridge itself, while leaving the semantic bridge
  intact. A separate certificate selects residual folds by membership;
  its projection to the core calculus and its repair of the minimal
  cascade are machine-checked. A fold-first saturation theorem and a
  domain-weighted flat compiler prove completeness of that new certificate
  for every admissible ground-closed trace.
\end{enumerate}

\paragraph{Isn't this just propositional logic?} A reader may object that
finite-domain many-valued logic reduces routinely to propositional logic:
ground out the quantifiers as finite folds and apply the companion paper's
completeness theorem. Three machine-checked facts say the reduction is not
routine. First, the naive lifted tableau is incomplete — crisp reflexivity
is valid and underivable (\cref{prop:naive}). Second, the \emph{bare}
grounding bridge is refuted: propositional consequence of the grounded
branch does not track finite-domain consequence, because propositional
valuations may assign the identity and equality atoms values no
finite-domain model induces (\cref{prop:barebridge}); repairing this needs
the rigid constraints of \cref{def:rigid}. Third, even with rigid
constraints, the repaired route initially proved completeness only for a
calculus with a propositional \emph{simulation macro} — a single rule
quoting the entire propositional tableau — which a proof-theoretic reading
wants eliminated; eliminating it is exactly the content of
\cref{thm:complete,thm:grounding}, and the documented failure of the
direct macro-elimination attempt (\cref{sec:simulation}) shows the
elimination has real content.

\paragraph{Method.} The development follows the project's rigor rules:
every claim is labeled, refutation is attempted before proof, and the Lean
artifact is the ground truth. The manuscript numbering is presentational;
the stable numbering and the full proofs live in the repository
documentation.

\section{Preliminaries: propositional NPL}\label{sec:prelim}

We summarize what is needed from the companion paper. The propositional
language has atoms and $\lnot,\wedge,\vee,\hmz$. The matrix $\FOUR$ is
$\{0,1\}^2$ with corners $\vT=(1,0)$, $\vF=(0,1)$, $\vB=(1,1)$,
$\vN=(0,0)$; negation swaps the two bits, $\wedge$ acts as $(\min,\max)$,
$\vee$ as $(\max,\min)$, and $\hmz$ as $(\min,\min)$ (bilattice consensus
\cite{ginsberg1988multivalued,fitting1991bilattices,arieli1996reasoning}).
The $\{\lnot,\wedge,\vee\}$-fragment is Belnap--Dunn first-degree
entailment \cite{belnap1977useful,dunn1976intuitive}. Four signs read the
bits: $\Tp$ ($t$-bit is $1$), $\Tm$ ($t$-bit is $0$), $\Fp$, $\Fm$
likewise on the $f$-bit. A four-signed analytic tableau calculus (sixteen
decomposition rules, two closure conditions) is sound and complete for
signed consequence over $\FOUR$, and — by the exact projection — over the
continuous semantics on $[0,1]^2$; signed tableaux over sets-of-values
signs are a standard paradigm \cite{haehnle1994automated}. All of this is
machine-checked in the same artifact this paper extends.

\section{Finite-domain quantified NPL}\label{sec:syntax}

\begin{definition}[Signature and syntax]\label{def:syntax}
Fix countable stocks of variables and predicate symbols and a signature
$\Sigma$ assigning each predicate symbol $P$ one arity
$\operatorname{ar}_{\Sigma}(P)\in\mathbb{N}$. Quantified
formulas are generated by
\[
\varphi \;::=\; P(x_0,\dots,x_k) \;\mid\; x = y \;\mid\; \lnot\varphi
\;\mid\; \varphi\wedge\varphi \;\mid\; \varphi\vee\varphi \;\mid\;
\varphi\hmz\varphi \;\mid\; \forall x\,\varphi \;\mid\; \exists x\,\varphi.
\]
A formula is $\Sigma$-well-formed iff every occurrence
$P(x_1,\ldots,x_k)$ has $k=\operatorname{ar}_{\Sigma}(P)$; the condition
is inherited structurally by compound formulas. The artifact retains a
total raw list syntax underneath this well-formed language.
[L:\,\lean{FiniteFO.QSignature}, \lean{FiniteFO.QFormula},
\lean{QFormula.WellFormed}]
\end{definition}

\begin{definition}[Models and assignments]\label{def:model}
For a fixed $n\in\mathbb{N}$, the domain is $\mathrm{Fin}(n{+}1)$ — finite
and nonempty by construction. A \emph{$\Sigma$-model} $M$ assigns to each
predicate symbol $P$ a map from tuples in
$\mathrm{Fin}(n{+}1)^{\operatorname{ar}_{\Sigma}(P)}$ to $\FOUR$.
Off-arity arguments are therefore absent from the model type. The artifact
retains a separate total raw list model for the tableau engine and connects
the two model classes by extension and restriction
(\cref{prop:sigbridge}). An \emph{assignment} $\rho$
maps variables to domain elements; $\rho[x{\mapsto}d]$ updates it.
[L:\,\lean{FiniteFO.QSigModel}, \lean{FiniteFO.QModel},
\lean{FiniteFO.Assignment}]
\end{definition}

\begin{definition}[Evaluation]\label{def:eval}
Evaluation extends the propositional $\FOUR$ clauses by crisp equality and
finite channel-wise quantifiers:
\[
V(x{=}y) =
\begin{cases}
\vT & \rho(x)=\rho(y),\\
\vF & \text{otherwise}.
\end{cases}
\]
\[
\begin{aligned}
V(\forall x\,\varphi)
  &= \Bigl(\textstyle\bigwedge_{d} t_d,\;\bigvee_{d} f_d\Bigr),\\
V(\exists x\,\varphi)
  &= \Bigl(\textstyle\bigvee_{d} t_d,\;\bigwedge_{d} f_d\Bigr).
\end{aligned}
\]
where $(t_d,f_d)$ is the value of $\varphi$ under $\rho[x{\mapsto}d]$ and
the extrema range over the finite domain. Signed satisfaction reads one
bit of the value, exactly as in the propositional case.
[L:\,\lean{FiniteFO.QSigModel.eval}, \lean{FiniteFO.qeval},
\lean{forallV4}, \lean{existsV4}]
\end{definition}

\begin{proposition}[Off-arity irrelevance]\label{lem:offarity}
If two models agree on every predicate tuple whose length is declared by
$\Sigma$, then they give the same value to every $\Sigma$-well-formed
formula under every assignment.
\end{proposition}

\begin{proof}
Structural induction on the formula. At a predicate atom,
well-formedness makes the model-agreement hypothesis applicable; equality
does not consult predicate data. The unary and binary cases follow from
the induction hypotheses and the corresponding deterministic $\FOUR$
operations. In a quantified case, apply the induction hypothesis at
every updated assignment and substitute the pointwise equalities into
the finite extrema. The hypothesis cannot be dropped: models may agree
on all declared unary tuples and disagree on the malformed atom $P()$.
[L:\,\lean{FiniteFO.qeval\_eq\_of\_agreeOn}]
\end{proof}

\begin{proposition}[Raw/signature semantic bridge]\label{prop:sigbridge}
Every $\Sigma$-model has a canonical total raw extension, obtained by
assigning an arbitrary fixed value to off-arity lists. Every raw model
restricts to a $\Sigma$-model. Restriction after extension recovers the
original $\Sigma$-model, while extension after restriction agrees with the
raw model on every declared-arity tuple. Consequently, for well-formed
$\Gamma,s\varphi$, semantic consequence over $\Sigma$-models is equivalent
to semantic consequence over raw models.
\begin{leanrefs}
[L:\,\lean{FiniteFO.QSigModel.toRaw\_restrict},\\
\quad\lean{FiniteFO.QModel.restrict\_toRaw\_agreeOn},\\
\quad\lean{FiniteFO.qconsequence4Sig\_iff\_qconsequence4}]
\end{leanrefs}
\end{proposition}

\begin{proof}
For a $\Sigma$-model, convert an arity-indexed tuple to its finite list and
back; finite-function extensionality and the list \texttt{ofFn/get}
identities recover the tuple. For a raw model, a declared-arity list is
recovered after conversion to the induced tuple, so the extension of its
restriction agrees on admitted tuples. For the consequence equivalence,
transfer branch and conclusion satisfaction along this agreement using
\cref{lem:offarity}. In the reverse direction, apply raw consequence
directly to the canonical extension of an arbitrary $\Sigma$-model.
\end{proof}

\begin{definition}[Branches; semantic consequence]\label{def:consequence}
A \emph{signed branch member} is a triple $s(\rho,\varphi)$ of a sign, an
assignment, and a formula — each member carries its \emph{own} assignment.
A \emph{branch} is a finite list of members; a model \emph{satisfies} a
branch iff it satisfies every member. \emph{Semantic consequence} is
domain-indexed: $\Gamma \qmod s\varphi$ iff every model over
$\mathrm{Fin}(n{+}1)$ that satisfies every member of $\Gamma$ satisfies
$s\varphi$.
\begin{leanrefs}
[L:\,\lean{FiniteFO.QSigned}, \lean{qsatSigned}, \lean{qsatBranch},\\
\quad\lean{FiniteFO.QSigModel.satSigned},
\lean{FiniteFO.QSigModel.satBranch},\\
\quad\lean{QConsequence4Sig}]
\end{leanrefs}
\end{definition}

\begin{proposition}[Quantifier duality]\label{prop:duality}
$\lnot\forall x\,\varphi$ and $\exists x\,\lnot\varphi$ have the same
value, as do $\lnot\exists x\,\varphi$ and $\forall x\,\lnot\varphi$:
channel-swap negation exchanges the two finite extrema. The nonempty
domain is used; over an empty domain the statement would depend on
conventions for vacuous quantifiers.
[L:\,\lean{FiniteFO.qeval\_neg\_all}, \lean{qeval\_neg\_ex}]
\end{proposition}

\begin{theorem}[Finite exact projection]\label{thm:projection}
Interpret predicates by continuous truth-objects in $[0,1]^2$ and evaluate
quantifiers by finite channel-wise infima/suprema. For every threshold
$0<\tau\le 1$, thresholding commutes with evaluation:
$\proj{\tau}(V_M(\varphi)) = V_{\proj{\tau}(M)}(\varphi)$, where
$\proj{\tau}(M)$ projects the predicate interpretations pointwise. Hence
the finite-domain metatheory, like the propositional one, can be developed
on $\FOUR$ and lifted. [L:\,\lean{FiniteFO.finite\_exact\_projection}]
\end{theorem}

\begin{remark}[Why finite domains]\label{rem:infinite}
Over an infinite domain the existential/supremum direction of
\cref{thm:projection} can fail when the supremum is not attained: every
instance can fall below the threshold while the supremum reaches it. The
finite, nonempty domain is therefore a load-bearing design choice, not a
convenience. (Recorded as a refutation attempt in the development.)
\end{remark}

\section{Grounding}\label{sec:groundtrans}

The finite domain makes every quantifier a finite fold, which suggests
translating quantified branches into propositional ones.

\begin{definition}[Ground atoms and the ground translation]\label{def:ground}
\emph{Ground atoms} are: two fixed identity atoms $\top_g$ and $\bot_g$
(fold units), a ground predicate atom for each predicate symbol and
argument list, and a ground equality atom for each pair of domain
elements. Ground atoms are encoded injectively into propositional atoms.
The translation $\grd_\rho$ maps a quantified formula under an assignment
to a propositional formula: predicate and equality atoms to their ground
atoms; connectives homomorphically; $\forall x\,\varphi$ to the finite
conjunction fold of $\{\grd_{\rho[x\mapsto d]}(\varphi)\}_d$ (with unit
$\top_g$) and $\exists$ to the disjunction fold (with unit $\bot_g$).
[L:\,\lean{FiniteFO.GroundAtom}, \lean{groundAtomCode},
\lean{FiniteFO.ground}, \lean{foldConj}, \lean{foldDisj}]
\end{definition}

\begin{lemma}[Grounding truth lemma]\label{lem:groundtruth}
For every model $M$ there is a propositional valuation $v_M$ (reading
ground atoms through $M$, with $\top_g\mapsto\vT$, $\bot_g\mapsto\vF$, and
equality atoms crisp) such that the propositional value of
$\grd_\rho(\varphi)$ under $v_M$ equals the quantified value of $\varphi$
under $M,\rho$. [L:\,\lean{FiniteFO.ground\_truth},
\lean{groundVal}]
\end{lemma}

\begin{definition}[Rigid constraints]\label{def:rigid}
The \emph{rigid constraints} form the finite propositional branch
containing $\Tp,\Fm$ for $\top_g$; $\Tm,\Fp$ for $\bot_g$; and for every
pair $(a,b)$ of domain elements, $\Tp,\Fm$ (respectively $\Tm,\Fp$) for
the ground equality atom of $(a,b)$ when $a=b$ (respectively $a\neq b$).
They pin exactly the part of a ground valuation that finite-domain models
do not choose freely; ground predicate atoms remain unconstrained.
[L:\,\lean{FiniteFO.rigidGroundConstraints}]
\end{definition}

\section{Calculi, and what fails}\label{sec:negative}

The propositional decomposition rules lift to signed quantified branch
members, acting on the member's own assignment. For self-containedness we
reproduce the sixteen propositional rows (from the companion paper;
``$A,B$'' adds both members to the branch, ``$A\mid B$'' splits it):

\begin{center}
\begin{tabular}{lcccc}
\toprule
 & $\Tp$ & $\Tm$ & $\Fp$ & $\Fm$ \\
\midrule
$\lnot\varphi$ & $\Fp\varphi$ & $\Fm\varphi$ & $\Tp\varphi$ & $\Tm\varphi$ \\
$\varphi\wedge\psi$ & $\Tp\varphi,\Tp\psi$ & $\Tm\varphi\mid\Tm\psi$
  & $\Fp\varphi\mid\Fp\psi$ & $\Fm\varphi,\Fm\psi$ \\
$\varphi\vee\psi$ & $\Tp\varphi\mid\Tp\psi$ & $\Tm\varphi,\Tm\psi$
  & $\Fp\varphi,\Fp\psi$ & $\Fm\varphi\mid\Fm\psi$ \\
$\varphi\hmz\psi$ & $\Tp\varphi,\Tp\psi$ & $\Tm\varphi\mid\Tm\psi$
  & $\Fp\varphi,\Fp\psi$ & $\Fm\varphi\mid\Fm\psi$ \\
\bottomrule
\end{tabular}
\end{center}

The eight finite quantifier rules come in two shapes, determined by which
channel bit the sign reads of the finite extrema:

\begin{center}\small
\begin{tabular}{ll}
\toprule
\emph{block} (add all $n{+}1$ instances, non-branching) &
$\Tp\forall,\ \Fm\forall,\ \Tm\exists,\ \Fp\exists$ \\
\emph{per-element} (one closed child per domain element) &
$\Tm\forall,\ \Fp\forall,\ \Tp\exists,\ \Fm\exists$ \\
\bottomrule
\end{tabular}
\end{center}

Closure, in the naive calculus, is an identical opposite-signed pair (same
assignment, same formula). [L:\,\lean{FiniteFO.QCloses}]

\begin{proposition}[The naive calculus is incomplete]\label{prop:naive}
Crisp reflexivity $\Tp(x{=}x)$ is valid over every finite domain, but the
empty branch does not derive it: no decomposition rule applies to an
equality atom and no closure pair exists. Machine-checked countermodel to
the completeness conjecture for \lean{QCloses}.
[L:\,\lean{FiniteFO.qeqRefl0\_valid},
\lean{qeqRefl0\_not\_derivable},
\lean{qcompleteness\_current\_refuted}]
\end{proposition}

\begin{definition}[Equality-completed calculus]\label{def:eqcalc}
Add four crisp-equality closure clauses: a branch closes if it contains
$\Tm(x{=}y)$ or $\Fp(x{=}y)$ under an assignment with $\rho(x)=\rho(y)$,
or $\Tp(x{=}y)$ or $\Fm(x{=}y)$ with $\rho(x)\neq\rho(y)$.
[L:\,\lean{FiniteFO.QClosesEq}]
\end{definition}

\begin{proposition}[The bare grounding bridge fails]\label{prop:barebridge}
Grounding alone does not reduce finite-domain consequence to propositional
consequence: $\Tp(x{=}x)$ is a finite-domain validity, but the grounded
sequent is \emph{not} a propositional consequence — a propositional
valuation may assign the ground equality atom of $(0,0)$ any value, while
every finite-domain model forces it to $\vT$. Machine-checked
countermodel; the rigid constraints of \cref{def:rigid} repair exactly
this gap. [L:\,\lean{FiniteFO.qeqRefl0\_ground\_not\_consequence4},
\lean{qeqRefl0\_ground\_branch\_not\_closes}]
\end{proposition}

\begin{example}[Why equality clauses are not enough]\label{ex:extensional}
Let $\rho(x)=\sigma(y)=d$ and consider the two-member branch
$\{\Tp(\rho,\,P\,x),\ \Tm(\sigma,\,P\,y)\}$. Every model gives both
members the same ground predicate value at $(P,[d])$, so the branch is
unsatisfiable; but the members are not an identical pair (different
formulas, different assignments), no rule decomposes literals, and no
equality clause applies. Closure must be allowed on \emph{equal
groundings}, not merely on identical members.
\end{example}

\section{The core extensional tableau}\label{sec:core}

\begin{definition}[Core calculus]\label{def:core}
The \emph{core extensional finite-domain tableau} extends the
equality-completed calculus of \cref{def:eqcalc} by two \emph{ground
closure clauses}: a branch closes if it contains members
$\Tp(\rho,\varphi)$ and $\Tm(\sigma,\psi)$ with
$\grd_\rho(\varphi)=\grd_\sigma(\psi)$, or $\Fp/\Fm$ likewise. All
sixteen propositional decomposition rules and all eight finite quantifier
rules are retained. The calculus contains \emph{no} propositional
simulation macro. Derivability: $\Gamma\qder s\varphi$ iff the branch
$\{\overline{s}\varphi\}\cup\Gamma$ closes, where $\overline{s}$ is the
opposite sign. [L:\,\lean{FiniteFO.QClosesExtCore},
\lean{QDerivesExtCore}]
\end{definition}

\begin{proposition}[Soundness]\label{prop:soundness}
If a branch closes in the core calculus, no finite model satisfies it;
hence $\Gamma\qder s\varphi$ implies $\Gamma\qmod s\varphi$.
[L:\,\lean{FiniteFO.QClosesExtCore.unsat}]
\end{proposition}

\begin{proof}[Proof sketch]
Induction on the closure derivation. Each closure condition is
unsatisfiable at its own members: an identical pair by sign
complementarity; an equality clause by crispness of $=$; a ground-closure
pair because \cref{lem:groundtruth} gives both members the same value
under any model, and opposite signs of one value cannot both hold. Each
decomposition rule preserves satisfiability downward into (one of) its
children by the corresponding local satisfaction equivalence — for the
quantifier rules, the signed reading of the finite extrema.
\end{proof}

\section{The completeness theorem}\label{sec:completeness}

\begin{theorem}[Semantic completeness of the core calculus]\label{thm:complete}
Fix $\Sigma$. Let $B$ be a finite $\Sigma$-well-formed branch of signed
quantified members over domain size $n{+}1$. If no model with domain
$\mathrm{Fin}(n{+}1)$ satisfies $B$, then $B$ closes in the core
calculus. With \cref{prop:soundness}, closure is \emph{exactly} finite
unsatisfiability. For every well-formed $\Gamma$ and $s\varphi$,
\[
\Gamma \qder s\varphi \quad\Longleftrightarrow\quad \Gamma \qmod s\varphi .
\]
The artifact proves the same equivalence first for its stronger raw
list syntax, proves the raw/signature semantic equivalence of
\cref{prop:sigbridge}, and composes the two equivalences.
\begin{leanrefs}
[L:\,\lean{FiniteFO.QClosesExtCore.complete\_of\_unsat},\\
\quad\lean{qclosesExtCore\_iff\_unsat},
\lean{QDerivesExtCore.complete},\\
\quad\lean{qDerivesExtCore\_iff\_qconsequence4Sig}]
\end{leanrefs}
\end{theorem}

\begin{proof}[Proof (engine; full details in the artifact)]
\emph{Step 1: the domain-weighted measure.} Define $\qsize$ on quantified
formulas by $\qsize=1$ on predicate and equality atoms, $+1$ for negation,
$+1$ over the sum for binary connectives, and
\[
\qsize(\forall x\,\varphi) \;=\; \qsize(\exists x\,\varphi) \;=\;
(n{+}1)\cdot\qsize(\varphi)+1 ,
\]
extended to branches by summation. Every rule of \cref{def:core} strictly
decreases the weight of the undecomposed segment: unary and binary rules
drop at least the head's outermost symbol; a block quantifier rule
replaces weight $(n{+}1)w+1$ by the $n{+}1$ instances of weight $w$ each;
a per-element rule replaces it by a single instance of weight
$w\le (n{+}1)w$. The case analysis is exhaustive over the eight formula
constructors and four signs.

\emph{Step 2: the engine.} By strong induction on the weight bound, a
branch split as $\mathit{todo} + \mathit{lits}$ (with $\mathit{lits}$
literals) that no finite model satisfies closes in the core: a literal
head moves to $\mathit{lits}$; a compound head is decomposed by its rule,
unsatisfiability transfers to each child by the corresponding local
satisfaction equivalence (the signed quantifier equivalences and the
propositional local-soundness table), the induction hypothesis applies at
strictly smaller weight, and monotonicity of closure under branch
extension [L:\,\lean{QClosesExtCore.mono}] restores the full
branch. Because the core's quantifier rules are natively $(n{+}1)$-ary,
no intermediate ``partially decomposed fold'' state ever arises
(contrast \cref{sec:simulation}).

\emph{Step 3: the literal stage.} For a branch of literals: if an equality
literal contradicts its crisp value, an equality clause closes; if two
literals carry $\Tp/\Tm$ (or $\Fp/\Fm$) with equal groundings, a ground
closure clause closes. Otherwise build the \emph{canonical model}: the
truth (falsity) bit of each ground predicate atom is set iff some $\Tp$-
($\Fp$-) literal grounds to it, read through the injective ground-atom
encoding. Positive literals are satisfied by their own witness; negative
literals because the opposite witness would be a ground closure pair,
excluded by assumption; equality literals by crispness, their four
unsatisfiable sign/value combinations being exactly the equality clauses.
This contradicts unsatisfiability, so a closure option applied.
\end{proof}

\begin{remark}[What the refutations bought]\label{rem:r5}
Three hypotheses of \cref{thm:complete} carry machine-checked failure
witnesses for their removal: without the ground closure clauses,
\cref{ex:extensional} is unsatisfiable but open; without the
$(n{+}1)$-weighting, decomposing a block quantifier \emph{increases} the
unweighted measure and the induction has no footing; without any one of
the four equality clauses, the corresponding literal configuration is
unsatisfiable but unmodelable and unclosable.
\end{remark}

\section{Constrained grounding reaches the core}\label{sec:grounding}

The artifact's first finite-domain completeness theorem ran through the
propositional calculus: if the grounded branch together with the rigid
constraints closes propositionally, an \emph{extensional macro calculus}
closes the quantified branch by a single simulation rule. That macro rule
is exactly what a proof-theoretic reading wants to eliminate.

\begin{theorem}[Constrained grounding, macro-free]\label{thm:grounding}
For every finite branch $B$: if the propositional branch
$\mathrm{Rigid}_n + \grd(B)$ closes in the propositional tableau
calculus, then $B$ closes in the core calculus of \cref{def:core}. No
simulation macro occurs: the propositional derivation is consumed through
propositional soundness, and the core closure is rebuilt by
\cref{thm:complete}. Concretely: a finite model of $B$ would induce, by
\cref{lem:groundtruth} and the rigid values, a propositional valuation
satisfying $\mathrm{Rigid}_n + \grd(B)$, contradicting propositional
soundness; hence $B$ is finitely unsatisfiable and
\cref{thm:complete} applies.
[L:\,\lean{FiniteFO.groundBranch\_closes\_to\_core},
\lean{satBranch\_groundVal\_rigid}]
\end{theorem}

Conversely, grounding is sound: the grounded branch of a satisfiable
quantified branch is satisfiable under the induced valuation
(\cref{lem:groundtruth}), so \cref{thm:grounding} composes with the
propositional completeness theorem of the companion paper into a
round-trip between the two levels.

\section{The simulation program: failure and repair}\label{sec:simulation}

An earlier phase of the project attempted \cref{thm:grounding} by
\emph{simulation}: induction on the propositional closure derivation,
maintaining a replay trace whose items include quantified members, rigid
atoms, and \emph{residual fold tails} — the partially decomposed remains
of a grounded quantifier, which are not groundings of any single
quantified formula. The layer produced a complete, machine-verified
dispatch analysis for propositional closure pairs (sixteen source
combinations per closure sign), a recursive ``tail consumer'' closing any
quantified formula against a fold tail it grounds into, and two further
machine-checked refutations of shortcut routes.

The program then hit a precise wall: when the propositional derivation
decomposes the grounding of a \emph{per-element-sign} quantifier
($\Tm\forall$ and its mirrors), the binary conjunction cascade produces a
right child asserting the pooled fold of the remaining instances in one
branch, while the core rule needs one closed child per domain element;
the pooled induction hypothesis cannot be redistributed, and
depth-preserving inversion does not make the induction well-founded. The
weighted-measure engine of \cref{thm:complete} dissolves the problem at
the root: the core's quantifier rules are $(n{+}1)$-ary, so the binary
cascade never appears. We record the dead end because the failure
analysis, not luck, produced the measure; the simulation layer remains in
the artifact as a verified study.

\begin{definition}[Replay traces and the two certificate predicates]
\label{def:replay}
Let $u=(\rho,\varphi)$ and write $q(S,u)$ for the quantified signed item
$(S,\rho,\varphi)$. A replay item is one of
\[
 q(S,u),\quad r(s),\quad C_S(L),\quad D_S(L),\quad
 C^q_S(U),\quad D^q_S(U),
\]
where $s$ is a ground signed formula, $L$ is a list of ground formulas,
and $U$ is a list of assignment--formula pairs. A replay trace is a finite
list of replay items. Its ground branch maps these six forms respectively to
the grounding of $q(S,u)$, $s$, the signed conjunction/disjunction folds of
$L$, and the signed folds of the groundings in $U$. Its quantified
projection keeps $q(S,u)$ and expands $C^q_S(U),D^q_S(U)$ to the list of
$q(S,u)$ for $u\in U$; the other forms contribute nothing.

A trace is \emph{admissible} iff every rigid item belongs to the printed rigid
constraint set, no unstructured $C_S(L)$ or $D_S(L)$ occurs, and the empty
structured folds $C^q_{\Tm}(\varnothing)$,
$C^q_{\Fp}(\varnothing)$, $D^q_{\Tp}(\varnothing)$, and
$D^q_{\Fm}(\varnothing)$ do not occur. These are exactly the four false
signed empty-fold cases.

The head-sensitive predicate $R(T)$ is the least proof predicate obtained by
applying every closure and decomposition scheme of \cref{def:core} to a
matching $q$-item in $T$, together with the corresponding rigid-equality
closure schemes, and the following fold schemes. For conjunction,
$S\in\{\Tp,\Fm\}$ consumes a nonempty list as head plus residual tail, while
$S\in\{\Tm,\Fp\}$ selects the list head; for disjunction the two sign pairs
are exchanged. Empty folds are erased precisely in the four admissible empty
cases. In all eight nonempty schemes the structured fold must be the
\emph{head} of $T$.

The membership-selecting predicate $R_{\in}(T)$ is the least predicate with
$R(T)\Rightarrow R_{\in}(T)$ and the following eight schemes, where the
displayed fold may occur anywhere in $T$:
\[
\begin{array}{c|c|c}
\text{fold} & \text{signs} & \text{recursive premise implying }R_{\in}(T)\\ \hline
C^q_S(u::U)\in T & S\in\{\Tp,\Fm\} &
 R_{\in}(q(S,u)::C^q_S(U)::T)\\
C^q_S(U)\in T,\ u\in U & S\in\{\Tm,\Fp\} &
 R_{\in}(q(S,u)::T)\\
D^q_S(u::U)\in T & S\in\{\Tm,\Fp\} &
 R_{\in}(q(S,u)::D^q_S(U)::T)\\
D^q_S(U)\in T,\ u\in U & S\in\{\Tp,\Fm\} &
 R_{\in}(q(S,u)::T)
\end{array}
\]
Each row abbreviates one constructor for each of its two signs. This definition
is a proof-object layer only: induction over its nine constructor schemes proves
$R_{\in}(T)\Rightarrow\textsc{QClosesExtCore}(q(T))$.
\begin{leanrefs}
[L:\,\lean{FiniteFO.ReplayItem.Admissible},
\lean{ReplayClosesCore}, \lean{ReplayClosesCoreMem},\newline
\lean{ReplayClosesCoreMem.toCore}]
\end{leanrefs}
\end{definition}

The old head-sensitive certificate conjecture also fails.  Over a two-element
domain, choose distinct ground predicate atoms $p_0,p_1$ and the admissible
trace consisting of a positive structured conjunction tail for
$[p_0,p_1]$ followed by the quantified source of $\Tm p_1$.  Its ground
branch is
\[
  \Tp\bigl(p_0\land(p_1\land\top_g)\bigr),\quad \Tm p_1,
\]
and closes by two positive-conjunction steps followed by atomic closure on
$p_1$.  Inverting a replay certificate first exposes $p_0$ ahead of the
residual one-element fold.  Since the exposed item is atomic, the residual
fold is no longer at the head where the certificate's fold constructors
require it, and the only possible atomic close would identify
$p_0=p_1$.  Injectivity of the ground-atom encoding contradicts that
equality.  Thus admissibility plus propositional closure does not imply an
existing replay certificate.  This is a failure of the head-sensitive
certificate representation, not of core soundness or completeness.
[L:\,\lean{FiniteFO.admissible\_ground\_replay\_bridge\_refuted}]

The repair is deliberately a new proof-object type rather than a revision
of the refuted claim.  The membership-selecting certificate embeds every
old certificate, but its structured-fold constructors may select a fold
occurrence anywhere in the trace.  All-item signs expose the head and
residual tail; branching signs select any represented item.  Induction on
the nine new constructor schemes proves sound projection into the
macro-free core calculus.  On the exact counterexample above, two
membership selections expose $p_1$ and an embedded atomic close finishes
the certificate.  Conversely, the earlier inadmissible empty-fold witness
still has no new certificate: sound projection would close the empty
quantified branch.  Thus the extension repairs the observed defect without
certifying the known false positive.
\begin{leanrefs}
[L:\,\lean{FiniteFO.ReplayClosesCoreMem.toCore},
\lean{FiniteFO.replayCascadeTrace\_\lb old\_\lb refuted\_\lb mem\_\lb verified},
\lean{FiniteFO.replayEmptyBadTailTrace\_\lb not\_\lb replayClosesCoreMem}]
\end{leanrefs}

We now prove the universal bridge for the repaired certificate.  For a
trace $T$ and quantified worklist $B$, let its flat normal form consist of
four signed atoms fixing the empty conjunction and disjunction folds, the
grounding of $B$, and exactly the rigid items retained by $T$.  A finite
reach relation generated by the membership-selecting constructors first
extends $T$ until every member of its quantified projection occurs as a
direct quantified item.  This removes the global-choice obstruction: all
branching alternatives may be materialized before compilation.

\begin{theorem}[Universal membership replay bridge]\label{thm:memreplay}
Every admissible replay trace whose ground branch closes has a
membership-selecting replay certificate.
\end{theorem}
\begin{proof}
Fold-first saturation gives an admissible extension $U$ containing the
original trace and directly representing every member of its quantified
projection.  Ground closure is monotone under this extension.  Any
valuation satisfying the flat normal form of $U$ satisfies its original
ground trace: direct quantified and rigid items transfer immediately;
structured folds transfer by induction over their represented finite
lists, with the four fixed atoms discharging empty residual folds.
Propositional soundness and completeness therefore close the flat normal
form of $U$.

It remains to compile that closure into the old head-sensitive certificate
over $U$.  We prove a stronger statement by induction on an upper bound for
the domain-weighted size of a worklist $\mathit{todo}$.  A second list
$\mathit{lits}$ contains only predicate or crisp-equality literals, and
every member of $\mathit{todo}\mathbin{+\!+}\mathit{lits}$ is required to
occur directly in $U$.  With empty $\mathit{todo}$, atomic closure
inversion exhausts the identity, quantified, and rigid sources and builds
an old certificate.  A literal at the head is moved to $\mathit{lits}$;
the resulting permutation preserves satisfaction and reduces the
worklist weight.

For negation and the three binary connectives, prepend the child item or
items prescribed by the old replay constructor.  The applicable
$\FOUR$ truth-table equation shows that satisfaction of each required
child flat branch implies satisfaction of the parent flat branch.  Parent
unsatisfiability and propositional completeness consequently close the
child branches.  Their measures are smaller: the outer connective unit is
removed, and a selecting branch also discards a subformula of positive
size.  Apply the induction hypothesis and the corresponding replay
constructor.

For an all-instance quantifier rule, prepend its $(n+1)$-element instance
block.  Its weight is $(n+1)\qsize(\varphi)$, one below the parent's
$(n+1)\qsize(\varphi)+1$, and satisfaction of all grounded instances
implies satisfaction of the parent finite fold.  For a selecting rule,
repeat the argument for each $d\in\operatorname{Fin}(n+1)$ using
membership of the selected grounding in the fold.  Positivity of $n+1$
gives
$\qsize(\varphi)\leq(n+1)\qsize(\varphi)<
(n+1)\qsize(\varphi)+1$; this includes $n=0$.  These cases cover both
quantifiers and all four signs.  The resulting old certificate for $U$
is transported backwards along the finite saturation path, yielding the
required membership-selecting certificate for $T$.
\end{proof}
\begin{leanrefs}
[L:\,\lean{FiniteFO.ReplayTrace.flat\_closes\_to\_replay},
\lean{FiniteFO.admissible\_ground\_replay\_bridge\_mem\_verified}]
\end{leanrefs}

The semantic route alone would still be insufficient: it produces core
closure rather than a value of the certificate type.  The theorem above
supplies the separate normalization and compilation argument.  It does not
resurrect the refuted old bridge, whose conclusion uses the strictly
smaller head-sensitive certificate.

\section{Related work}\label{sec:related}

\emph{Finitely-valued first-order tableaux.} Complete tableau calculi for
arbitrary finite many-valued first-order logics are classical:
Carnielli's systematization \cite{carnielli1987systematization} proves
abstract completeness, model existence, compactness and
L\"owenheim--Skolem theorems for tableaux with \emph{distribution
quantifiers}, and H\"ahnle's monograph \cite{haehnle1994automated}
generalizes the sign language to sets-as-signs. We therefore claim no
novelty for many-valued first-order tableau completeness as such, and
NPL's four threshold signs are themselves sets-as-signs instances
($\{\vT,\vB\}$, $\{\vF,\vN\}$, $\{\vB,\vF\}$, $\{\vT,\vN\}$). The present
calculus differs in format and scope: quantifiers range over a
\emph{fixed finite nonempty domain} and are decomposed by $(n{+}1)$-ary
instance rules rather than through value distributions over arbitrary
domains; the closure conditions include crisp-equality clauses and
ground-closure (extensional) clauses, neither of which is supplied by the
generic per-formula calculi cited above; and the completeness proof
is driven by the domain-weighted measure, which has no counterpart in the
distribution-quantifier setting. Two delimitations are worth making
explicit. Semantically, fixing the domain \emph{in the judgment} is a
design choice, not a lack: the intended applications quantify over known
finite data, and the domain-indexed consequence relation
(\cref{def:consequence}) is the natural primitive there — the
distribution-quantifier tradition answers a different question
(consequence over all structures). Proof-theoretically, one may ask
whether the core calculus is an instance of the sets-as-signs framework
whose generic completeness argument applies off the shelf: the
decomposition rules are, but the crisp-equality clauses and the
ground-closure clauses act on \emph{pairs of members across different
assignments} through the grounding translation, which is outside the
per-formula format of the generic argument; \cref{ex:extensional} is
precisely a branch that every sets-as-signs closure notion keyed to
identical signed formulas leaves open.

\emph{Quantified Belnap--Dunn logics.} First-order FDE and its tableaux
are textbook material \cite[Ch.~22]{priest2008introduction}, on the classical
foundations of \cite{belnap1977useful,dunn1976intuitive}. Recent work
studies free quantification over four-valued and fuzzy bilattice-valued
logics with dual-domain semantics for non-denoting terms
\cite{behounek2023free}; that axis (free logic, unrestricted domains, no
mechanized proof theory) is orthogonal to ours (total denotation, fixed
finite domains, machine-checked exact characterization). The
harmonization connective $\hmz$ in the object language and the threshold
reading remain the NPL-specific propositional base, positioned against
the bilattice literature in the companion paper.

\emph{Machine-checked completeness.} Formalized metatheory exists for
classical first-order proof systems, including an Isabelle/HOL coinductive
framework instantiated for Gentzen and tableau variants of classical logic
\cite{blanchette2017soundness}. TableauxRocq gives a recent Rocq verification
of soundness and certificate checking for a classical first-order tableau
\cite{rosain2026tableauxrocq}; it does not claim a many-valued completeness
result. Formalized paraconsistent metatheory also predates this work
\cite{villadsen2017formalizing}. Our updated literature pass (2026-08-16) found no
machine-checked completeness theorem for a many-valued quantified tableau
calculus. This supports only the search-scoped conclusion that, to our
knowledge, the present development appears to be the first.

\section{Open problems}\label{sec:open}

(1)~\emph{Infinite domains}: the projection argument fails at unattained
suprema (\cref{rem:infinite}); a quantified NPL over infinite domains
needs either attained-extrema semantics or a genuinely different bridge.
(2)~\emph{Function symbols and non-crisp equality}: the present language
is function-free and its equality never takes glut or gap values; both
restrictions are design choices worth revisiting.
(3)~\emph{Complexity}: finite-domain consequence is decidable by
finiteness of the model space; its exact complexity as a function of
domain size and formula size is open here.
(4)~\emph{Certificate complexity and normalization}: the repaired
membership-selecting certificate is complete by \cref{thm:memreplay}, but
the present proof materializes all quantified-projection members before
compilation.  Tight upper bounds on certificate growth, and a focused or
zipper-style compiler that avoids unnecessary materialization, remain
open.

\appendix
\section{The Lean artifact}\label{app:lean}
The results of this paper are verified in the \lean{Nullivance} Lean~4
library, principally module \lean{Nullivance.FiniteFO}, and share the
artifact of the companion paper. \textbf{Availability.} The complete source
artifact is supplied with this release, including citation/license
metadata, exact dependency locks, both manuscripts, and verification scripts.
The tagged source archive is available from
\href{https://github.com/lamtung0487-droid/nullivance-logic/releases/tag/v0.7.1}{GitHub release v0.7.1}
and archived under DOI
\href{https://doi.org/10.5281/zenodo.21964600}{10.5281/zenodo.21964600}.
\textbf{Toolchain.}
Lean~4 (\texttt{leanprover/lean4:v4.32.1}) and mathlib \texttt{v4.32.1} at
commit \texttt{520045ab14e26149ee970e2e617ca04b09bde5d6}; run
\texttt{scripts/Verify-Release.ps1}. \textbf{Coverage policy.} The build is
\texttt{sorry}-free. The headline theorems are stated directly against
\lean{QClosesExtCore}; replay-certificate results use the separately defined
proof-object predicates of \cref{def:replay} and project into that core.
Axiom audit: every named declaration depends only on
\lean{propext}, \lean{Classical.choice}, \lean{Quot.sound}.
\textbf{Key declarations.}
\begin{leanrefs}
Semantics: \lean{QSigModel}, \lean{qeval}, \lean{qsatBranch}.\\
Projection: \lean{finite\_exact\_projection}.\\
Refutations: \lean{qcompleteness\_current\_refuted},
\lean{qeqRefl0\_not\_derivable}.\\
Calculus: \lean{QClosesExtCore};
soundness: \lean{QClosesExtCore.unsat}.\\
Completeness engine: \lean{qsize}, \lean{qweightB},
\lean{qclosesCore\_todo}, \lean{qclosesCore\_lits}.\\
Headline: \lean{QClosesExtCore.complete\_of\_unsat},\\
\quad\lean{qDerivesExtCore\_iff\_qconsequence4Sig}.\\
Grounding: \lean{ground\_truth},
\lean{groundBranch\_closes\_to\_core}.\\
Replay repair: \lean{ReplayClosesCoreMem},
\lean{ReplayClosesCoreMem.toCore},
\lean{replayCascadeTrace\_\lb old\_\lb refuted\_\lb mem\_\lb verified}.\\
Replay completeness: \lean{ReplayTrace.flat\_\lb closes\_\lb to\_\lb replay},
\lean{admissible\_\lb ground\_\lb replay\_\lb bridge\_\lb mem\_\lb verified}.
\end{leanrefs}

{\footnotesize
\bibliographystyle{plain}
\bibliography{../../references/bibliography}
}

\end{document}
